% DRAFT of the method section of the intervention-transfer study for the course report (written 2026-10-26 before any outcome of the study existed; the numbers are those of freeze.json).
% Not compiled into the report yet: it is pasted into main.tex when the report is rewritten around the study. Every claim here is a design fact, not a result.
\section{Does an intervention keep its effect under quantization?}
\label{sec:transfer}

\paragraph{Intervention and cells.}
We remove the full-precision decoder contribution $v$ of one SAE feature at the positions where the feature is active in full precision (FP), for the layer-12 residual of Gemma-2-2B. The edit and the readout are computed in float32; every cell subtracts its own unedited baseline, and a zero edit must give exactly zero. Let $D[p,s]$ be the mean over the edited positions of the centered change of the output logits, where the prefix state $p\in\{\mathrm{FP},\mathrm{Q}\}$ is the residual (and hence the SAE activation) \emph{before} the edit and the suffix state $s\in\{\mathrm{FP},\mathrm{Q}\}$ is the set of weights \emph{after} it. The context term $C=D[\mathrm{Q},\mathrm{FP}]-D[\mathrm{FP},\mathrm{FP}]$, the weight term $W=D[\mathrm{FP},\mathrm{Q}]-D[\mathrm{FP},\mathrm{FP}]$ and the interaction $I=D[\mathrm{Q},\mathrm{Q}]-D[\mathrm{Q},\mathrm{FP}]-D[\mathrm{FP},\mathrm{Q}]+D[\mathrm{FP},\mathrm{FP}]$ satisfy $D[\mathrm{Q},\mathrm{Q}]-D[\mathrm{FP},\mathrm{FP}]=C+W+I$ exactly. A fifth cell edits the quantized state with its \emph{own} activation of the feature (native amplitude). Weights are quantized by round-to-nearest per output channel at 6 and 4 bits; 16 bits is the rounding-only reference.

\paragraph{Floor, eligibility, controls.}
Because the model is executed in bfloat16, even 16-bit weights change the logits. Per direction kind we fix a numerical floor as the 95th percentile of the 16-bit error norm over \emph{calibration} features, and call an item eligible if its full-precision effect norm exceeds five floors. The normalized transfer error is $\lVert D[\mathrm{Q},\mathrm{Q}]-D[\mathrm{FP},\mathrm{FP}]\rVert/(\lVert D[\mathrm{FP},\mathrm{FP}]\rVert+\text{floor})$. Each edit is also made along a matched swapped direction (another feature's decoder row) and three random norm-matched directions orthogonal to the real one, so that what is specific to the feature can be told from what any direction would do.

\paragraph{Populations.}
The feature level uses 404 language-selective features (up to 16 contexts each, the dev sentences of largest activation; 6,464 feature--context pairs, 32,320 items per state), split by a hash of the feature id into calibration and held-out halves. The group level ablates each language's 64-feature group and three matched control groups on 27 languages $\times$ 300 FLORES sentences and decomposes the net harm (group minus controls). An external panel of 1,800 Wikipedia passages in six languages (revision pinned, text hashed) repeats both levels on fresh text with its own contexts.

\paragraph{Pre-specified analysis.}
All intervals are 2,000-replicate crossed bootstraps over features and source-sentence clusters (seed 20261026). Before any outcome we fixed: the expectation that $|C|>|W|$ at the group level (A3, the loss-level split of \citet{chang2025} suggests it, so it is not a contribution by itself); H1, that more than half of the error remains after the best rescaling of the full-precision effect and after using the native amplitude; H2, the descriptive reproducibility of $C$, $W$, $I$ between the calibration and the held-out half; H3 and the utility rule: a ridge model adding the sensitivity to the \emph{other} directions of the same context to an activation baseline must reduce the held-out absolute error by at least 10\% against the strongest baseline (activation, or activation plus a gradient attribution), with a paired interval above zero, at 6 and 4 bits separately; H4, that the averaged fingerprint of a feature transfers differently from its single contexts. A secondary question asks whether the Pearson survival of the feature adds to these predictors.
